% Context from chapters/wavelets.tex; for mathematical reference, not compilation.
\section{2-D Wavelets}

                                                   
\subsection{Anisotropic Wavelets}

The anisotropic basis on $\TT^2$ is the tensor product of the complete one-dimensional periodized basis $\mathcal B$ (including its coarsest scaling functions):
\eq{\{b_1\otimes b_2:b_1,b_2\in\mathcal B\}.}
Its wavelet--wavelet elements have the form
\eql{\label{eq-aniso-wav}
 \psi_{(j_1,j_2),(n_1,n_2)}(x_1,x_2)
 =\psi_{j_1,n_1}^{P}(x_1)\psi_{j_2,n_2}^{P}(x_2).
}
The basis also includes wavelet--scaling, scaling--wavelet, and scaling--scaling products.
The anisotropic 2-D transform follows the same separable construction as the 2-D FFT in Section~\ref{sec-fft-2d}. Treat the image $a_J \in \RR^{2^{-J} \times 2^{-J}}$ as a matrix, apply the 1-D FWT to each row, and then apply it to each column. The total cost is $O(N)$ for $N=2^{-2J}$ pixels.

\myfigure{
\tabtrois{
\image{wavelets}{.3}{wavalgo-2d-aniso-0}&
\image{wavelets}{.3}{wavalgo-2d-aniso-1}&
\image{wavelets}{.36}{wavalgo-2d-aniso-2}\\
Image $f$ & Row transform & Column transform.
}
}{ 
	Steps of the anisotropic wavelet transform.  
}{fig-wavalgo-2d-aniso}

                                                   
\subsection{Isotropic Wavelets}\label{sec-isotropic-wav}

The anisotropic atoms~\eqref{eq-aniso-wav} have independent scales in the two coordinate directions. For compactly supported wavelets, their supports lie in axis-aligned rectangles of size proportional to $2^{j_1}\times2^{j_2}$. This directional bias can produce visible artifacts in denoising and compression when it does not match the image geometry.

\myfigure{
\image{wavelets}{.3}{wavcoefs-iso-aniso-2}
\image{wavelets}{.3}{wavecoefs-iso-aniso-1}
}{ 
	Anisotropic (left) versus isotropic (right) wavelet coefficients.  
}{fig-wavecoefs-iso}

An alternative uses a common scale in both directions. Tensor products of one-dimensional approximation spaces give a two-dimensional multiresolution
\eq{
	L^2(\RR^2) \supset \ldots \supset V_{j-1} \otimes V_{j-1} \supset V_j \otimes V_j \supset V_{j+1} \otimes V_{j+1} \supset \ldots \supset \{0\}.
}
We write
\eq{
	V_j^O \eqdef V_j \otimes V_j
} 
for this isotropic 2-D approximation space.

Recall that the tensor product of two closed subspaces $V_1,V_2\subset L^2(\RR)$ is
\eq{
	V_1 \otimes V_2 = \text{Closure}\pa{ \Span\enscond{f_1(x_1)f_2(x_2) \in L^2(\RR^2)}{ f_1 \in V_1, f_2 \in V_2 } }.
}
If $(\phi_k^s)_k$ are orthonormal bases of $V_s$, their tensor products $(\phi_{k_1}^1(x_1)\phi_{k_2}^2(x_2))_{k_1,k_2}$ form an orthonormal basis of $V_1 \otimes V_2$.

The tensor product distributes over these orthogonal sums
\eq{
		( V_j \oplus^\bot W_j) \otimes ( V_j \oplus^\bot W_j)
		= 
		V_j^O \oplus^\bot W_j^V\oplus^\bot W_j^H\oplus^\bot W_j^D
		\qwhereq
		\choice{
		W_j^V \eqdef (V_j \otimes W_j), \\
		W_j^H \eqdef (W_j \otimes V_j),  \\
		W_j^D \eqdef (W_j \otimes W_j).
		}
}
The labels $\{V,H,D\}$ stand for \textit{Vertical, Horizontal, Diagonal} detail spaces.
 
The following diagram summarizes the nested spaces and their detail components:
\begin{center}
	\image{wavelets}{.8}{embeded-spaces-2d}
\end{center}

For $j \i